\section{研究结论、局限与下一步问题}
\label{sec:discussion}

\subsection{可以确认的系统贡献}
本项目将大模型决策与数学计算分开实现：模型负责形成可检验的研究动作，求解器负责计算，独立评测负责检验。公开决策说明、实际模型原始输出、工具结果、阶段记忆和审阅意见可相互核对。这一结构使失败能够成为下一步输入，也使报告中的数值能够追溯到逐例计算。它体现了一个可运行科研智能体的开发与应用，而不是将预先写好的研究结论排成对话。

同时，系统不应被称为完全自主的科学家。问题选择、算子库、合成分布、工具接口、预算和统计协议由开发者设计；智能体在这个受限研究空间中组织假设、组合组件与修订方法。数学证明需要人工与符号层面的审查；独立critic调用仍来自同一配置模型，不能替代独立研究团队复现。两轮相继修订的真实轨迹也不足以证明这种规划策略优于随机、专家或其它模型；这需要额外的等预算、同问题、多次研究对照。

\subsection{数学与统计上的边界}
观测自加权偏差命题针对简化的标量对数正态模型。坐标误差、加性噪声、多锚点几何、截尾与局部优化同时出现时，其具体偏差并不由同一个标量公式决定。对数残差因此是一项受动机支持的候选，而非通用最优解。

镜像反例表明，数据本身的全局不可辨识性不能靠换损失函数消除。局部Jacobian满秩与求解器收敛，均不能排除另一处远离当前解的等价轨迹。局部误差界需要真值邻域、距离下界和曲率条件；它不能直接作为每个合成案例的保证书。

交叉验证衡量观测窗口内的预测适配，不保证选出使当前位置误差最小的运动类别。异常验证观测、非恒速运动和过少样本均可能破坏选择。共形区域的边际覆盖依赖完整案例之间的交换性，压力分布与拒绝后的条件子集超出了该保证。

\subsection{后续可检验的新问题}
本项目进一步提出一个适合继续研究的问题：在给定污染比例上界、最低几何信息量和有限求解预算下，能否构造允许拒绝的当前位置估计器，同时对静止情形给出非劣界、对移动情形给出误差改善，并控制拒绝率？这一问题必须明确分布类，否则镜像歧义会给出直接的不可能性障碍。

下一轮工作应首先处理可识别的部分：引入少量破除镜像的方向或地图信息，研究其最小数量；对非视距与设备标定漂移建立明确污染模型；用真正按设备或轨迹划分的现场数据检验区域覆盖；最后对不同智能体规划器开展多次独立重复与预算匹配的对照。它们是本项目留下的研究任务，而非本次已经完成的效果声明。
